The performance of the ANNs is assessed using the receiver operating characteristic (ROC) curve. To evaluate model generalization and check for overfitting, a 5-fold cross-validation strategy is used. This approach provides five independent test sets, offering a robust estimate of the model's performance without introducing bias from reliance on a single test set. The ROC curves from the 5-fold validation are shown in Figure~\ref{fig:vh_wh_zdom_roc} for the \zdom{} training and in Figure~\ref{fig:vh_wh_zdep_roc} for the \zdep{} training.

The ANN score distributions across the five folds are presented in Figure~\ref{fig:vh_wh_zdom_kfold_scores} (\zdom{}) and Figure~\ref{fig:vh_wh_zdep_kfold_scores} (\zdep{}). These distributions exhibit strong agreement between training and test samples, with consistent shapes across all folds, indicating no evidence of overfitting. However, there is an exception to this in Figure~\ref{fig:vh_wh_wz_zdep_kfold_scores}, where there is some disagreement seen in these plots. This is due to the limited statistics of the training sample, and would disappear with more MC events.

Based on these results, the ANNs are retrained using the full training dataset, taking advantage of all available information.
